% !TeX root=../main.tex
\chapter{مفاهیم پایه}
\label{ch:background}

\section{مقدمه}

طراحی یک برنامه برای کودکان زیر سه سال نیازمند شناخت رشد ارتباط و زبان، نقش بازی، اهمیت تعامل والد و محدودیت‌های رسانهٔ دیجیتال است. فعالیت جذاب به‌تنهایی «گفتاردرمانی» محسوب نمی‌شود و یک برنامهٔ آموزشی نیز نباید جای ارزیابی تخصصی را بگیرد. در این فصل مفاهیم اصلی رشد گفتار و زبان و ارتباط آن‌ها با قابلیت‌های نوواژه به‌صورت فشرده بررسی می‌شوند.

\section{ارتباط، گفتار و زبان}

ارتباط، گفتار و زبان مفاهیمی مرتبط اما متفاوت‌اند. \textbf{ارتباط} گسترده‌ترین مفهوم است و هر روشی را در بر می‌گیرد که کودک با آن نیاز، احساس یا قصد خود را منتقل می‌کند. گریه، لبخند، نگاه، اشاره، نزدیک‌کردن یک شیء به والد، تولید آوا و بیان واژه همگی می‌توانند رفتار ارتباطی باشند. در سال نخست، بسیاری از تعامل‌ها پیش‌زبانی‌اند؛ برای مثال کودک با نگاه میان والد و یک اسباب‌بازی، زمینهٔ توجه مشترک را ایجاد می‌کند. اگر والد در همان لحظه نام شیء را بیان کند، صدا، معنا و تجربهٔ اجتماعی به یکدیگر پیوند می‌خورند.

\textbf{گفتار} یکی از روش‌های بیان زبان و حاصل هماهنگی تنفس، حنجره، زبان، لب‌ها و دیگر اندام‌های گفتاری است. دقت تولید صداها، روانی و کیفیت صوت از جنبه‌های گفتار به شمار می‌روند \cite{NIDCDMilestones2022}. ممکن است کودکی معنای یک واژه را بداند، اما هنوز برخی صداهای آن را دقیق تلفظ نکند. بنابراین شناخت یک تصویر پس از شنیدن نام آن با توانایی تولید درست همان واژه یکسان نیست.

\textbf{زبان} سامانه‌ای از نشانه‌ها و قواعد برای درک و انتقال معناست. زبان دریافتی به فهم پیام دیگران، مانند دنبال‌کردن دستور ساده یا پیداکردن یک تصویر پس از شنیدن نام آن، مربوط می‌شود. زبان بیانی نیز توانایی انتقال معنا با حرکت، آوا، واژه یا ترکیب واژه‌هاست. در سال‌های آغازین معمولاً کودک واژه‌های بیشتری را می‌فهمد تا بیان کند. به همین دلیل تمرین انتخاب تصویر می‌تواند بخشی از واژگان دریافتی را تقویت کند، اما دربارهٔ تلفظ، جمله‌سازی یا استفادهٔ خودانگیختهٔ واژه اطلاعات کاملی نمی‌دهد.

یادگیری واژه تنها حفظ یک برچسب نیست. کودک باید شکل شنیداری واژه را با معنای آن پیوند دهد، نمونه‌های متفاوت آن را تشخیص دهد و واژه را در موقعیت مناسب به کار ببرد. برای مثال، شناخت تصویر یک سیب نخستین گام است؛ تعمیم زمانی رخ می‌دهد که کودک سیب واقعی با رنگ یا اندازهٔ متفاوت را نیز بشناسد یا برای درخواست آن از اشاره و واژه استفاده کند. ازاین‌رو فعالیت دیجیتال باید با گفت‌وگو و اشیای واقعی ادامه پیدا کند.

\section{رشد گفتار و زبان تا سه‌سالگی}

سه سال نخست زندگی دوره‌ای سریع برای رشد ارتباط و زبان است. محیطی که در آن کودک صدای دیگران را می‌شنود، پاسخ مناسب دریافت می‌کند و فرصت بازی و گفت‌وگو دارد، از این رشد حمایت می‌کند \cite{NIDCDMilestones2022}. بااین‌حال، سرعت رشد در کودکان یکسان نیست و نقاط عطف باید به‌عنوان راهنمای مشاهده استفاده شوند، نه معیار تشخیص بر اساس یک رفتار منفرد. شنوایی، سابقهٔ رشد، زبان‌های خانواده و کیفیت تعامل‌های روزانه نیز باید در تفسیر توانایی کودک در نظر گرفته شوند \cite{ASHAMilestones2024}.

\begin{table}[ht]
\centering
\caption{نمونه‌هایی از مسیر عمومی رشد ارتباط و زبان تا سه‌سالگی \cite{NIDCDMilestones2022,ASHAMilestones2024}}
\label{tab:language-milestones}
\renewcommand{\arraystretch}{1.25}
\begin{tabular}{|p{0.15\textwidth}|p{0.34\textwidth}|p{0.37\textwidth}|}
\hline
\textbf{بازهٔ سنی} & \textbf{درک و ارتباط} & \textbf{تولید آوا و زبان} \\
\hline
تولد تا ۶ ماه & واکنش به صدا، توجه به چهره و آرام‌شدن با صدای آشنا & گریه‌های متفاوت، لبخند، خنده و بازی با آواها \\
\hline
۶ تا ۱۲ ماه & توجه مشترک، شرکت در بازی‌هایی مانند دالی و درک برخی نام‌های آشنا & زنجیره‌های آوایی، تقلید صدا، اشاره و ظهور نخستین واژه‌ها \\
\hline
۱۲ تا ۲۴ ماه & دنبال‌کردن دستور ساده و اشاره به شیء یا تصویر پس از شنیدن نام & افزایش واژه‌ها و نزدیک‌شدن به ترکیب دو واژه \\
\hline
۲۴ تا ۳۶ ماه & درک پرسش‌های روزمره و شناخت نام بسیاری از اشیا & عبارت‌های کوتاه، نامیدن و درخواست‌کردن با واژه \\
\hline
\end{tabular}
\end{table}

توجه مشترک یکی از پایه‌های یادگیری واژه است. هنگامی که کودک و والد هم‌زمان به یک شیء توجه دارند و والد نام آن را بیان می‌کند، احتمال شکل‌گیری پیوند میان واژه و مرجع افزایش می‌یابد. نمایش هم‌زمان تصویر و پخش صدا در نوواژه نیز همین پیوند دیداری و شنیداری را هدف می‌گیرد، اما همراهی والد آن را به تجربه‌ای اجتماعی تبدیل می‌کند.

تکرار برای تثبیت واژه ضروری است، ولی تکرار باید با تنوع همراه باشد. شنیدن چندبارهٔ نام «ماشین» مفید است، اما کودک باید آن را برای ماشین اسباب‌بازی، تصویرهای متفاوت و وسیلهٔ واقعی در خیابان نیز بشنود. برنامه می‌تواند تکرار را منظم کند؛ تعمیم آموخته به زندگی روزمره با مشارکت خانواده صورت می‌گیرد.

\section{تأخیر زبانی و گفتاردرمانی}

تأخیر زبانی زمانی مطرح می‌شود که مهارت‌های زبانی کودک آهسته‌تر از مسیر مورد انتظار ظاهر شوند. بعضی کودکان بیشتر در بیان واژه مشکل دارند و برخی در درک و بیان هر دو با دشواری روبه‌رو هستند. اختلال گفتار ممکن است با تولید نادرست صداها، ناروانی یا دشواری حرکتی همراه باشد و اختلال زبان می‌تواند واژگان، درک معنا یا ساخت جمله را درگیر کند \cite{NIDCDMilestones2022,ASHALateLanguage}.

یک تفاوت سنی به‌تنهایی علت یا نوع مشکل را مشخص نمی‌کند. وضعیت شنوایی، ویژگی‌های رشدی، فرصت تعامل و زبان یا زبان‌های خانواده بر مسیر ارتباط اثر دارند. چندزبانگی به‌خودی‌خود اختلال ایجاد نمی‌کند و ارزیابی کودک چندزبانه باید همهٔ زبان‌های او را در نظر بگیرد. در صورت نگرانی دربارهٔ از‌دست‌رفتن مهارت‌ها، واکنش‌ندادن به صدا یا نام، استفاده‌نکردن از حرکت‌های ارتباطی، دشواری پایدار در درک یا پیشرفت بسیار کم، مراجعه به متخصص ضروری است.

آسیب‌شناس گفتار و زبان (\lr{Speech-Language Pathologist}) برای ارزیابی تنها تعداد واژه‌ها را نمی‌شمارد. گزارش والد، مشاهدهٔ بازی و تعامل، درک دستورها، استفاده از اشاره، نمونهٔ زبان، تولید صدا و وضعیت اندام‌های گفتاری می‌توانند بخشی از بررسی باشند. ارزیابی شنوایی نیز مهم است. نتیجهٔ این فرایند مشخص می‌کند که کودک به پایش، آموزش خانواده، درمان مستقیم یا ارجاع تکمیلی نیاز دارد.

مداخلهٔ زودهنگام با هدف استفاده از انعطاف‌پذیری سال‌های نخست و تقویت فرصت‌های یادگیری در محیط طبیعی انجام می‌شود \cite{ASHAEarlyIntervention}. در کودکان خردسال، والد معمولاً بخش اصلی مداخله است؛ زیرا بیشترین فرصت تمرین در غذاخوردن، حمام، لباس‌پوشیدن، کتاب‌خواندن و بازی رخ می‌دهد. پژوهش‌ها نیز نشان داده‌اند آموزش راهبردهای زبانی به والد می‌تواند بر مهارت‌های ارتباطی کودک اثر مثبت داشته باشد \cite{RobertsKaiser2011,RobertsCurtis2019}.

راهبردهای رایج حمایت از زبان عبارت‌اند از:

\begin{itemize}
\item دنبال‌کردن علاقهٔ کودک و نامیدن چیزی که به آن نگاه می‌کند؛
\item پاسخ‌دادن سریع به نگاه، اشاره، آوا یا واژه؛
\item گسترش بیان کودک؛ برای مثال تبدیل «توپ» به «توپ قرمز»؛
\item ایجاد مکث و فرصت برای پاسخ، بدون فشار برای تکرار؛
\item استفاده از جمله‌های کوتاه و طبیعی همراه با حرکت و شیء واقعی؛
\item تکرار واژه در فعالیت‌ها و موقعیت‌های متفاوت.
\end{itemize}

نوواژه می‌تواند تصویر، صوت، تمرین کوتاه و سابقهٔ پاسخ را در اختیار خانواده بگذارد، اما تشخیص یا درمان پزشکی نیست. برنامه تلفظ، شنوایی، تعامل اجتماعی و زبان خودانگیخته را به‌طور جامع ارزیابی نمی‌کند و نتیجهٔ بازی نباید به‌عنوان گزارش بالینی تفسیر شود.

\section{اثر بازی بر یادگیری زبان}

بازی یکی از طبیعی‌ترین موقعیت‌های یادگیری در کودکی است. کودک در بازی به شیء توجه می‌کند، نقش می‌سازد، نوبت می‌گیرد و برای ادامهٔ تعامل از حرکت، آوا و واژه استفاده می‌کند. پژوهش‌ها میان کیفیت بازی و رشد زبان رابطه نشان داده‌اند، هرچند نوع بازی، سن کودک و مشارکت بزرگسال بر نتیجه اثر می‌گذارند \cite{Toub2018PlayLanguage}.

یادگیری مبتنی بر بازی (\lr{Game-Based Learning}) به استفاده از یک فعالیت دارای هدف و قواعد برای یادگیری اشاره دارد. بازی‌وارسازی (\lr{Gamification}) نیز افزودن عناصری مانند امتیاز، مرحله یا بازخورد به یک فعالیت آموزشی است. مرور پژوهش‌های آموزش کودکی نشان می‌دهد بازی می‌تواند انگیزه و مشارکت را افزایش دهد، ولی این نتیجه به طراحی فعالیت و هدف آموزشی وابسته است \cite{Alotaibi2024GameBased}.

یک بازی زبانی مناسب می‌تواند از چند مسیر اثرگذار باشد:

\begin{enumerate}
\item مرحلهٔ کوتاه و هدف روشن، توجه کودک را حفظ می‌کند؛
\item تکرار واژه در پرسش‌های متفاوت، فرصت تمرین ایجاد می‌کند؛
\item انتخاب تصویر پس از شنیدن واژه، بازیابی فعال را درگیر می‌کند؛
\item بازخورد فوری رابطهٔ میان انتخاب و نتیجه را روشن می‌سازد؛
\item مشارکت والد هر تصویر را به پرسش، اشاره و نمونهٔ واقعی تبدیل می‌کند.
\end{enumerate}

بازخورد باید کوتاه، مثبت و مرتبط با هدف باشد. پویانمایی طولانی، صدای شدید یا پاداش نامرتبط ممکن است توجه را از واژه دور کند. پژوهش تعامل لمسی نیز نشان داده است ویژگی‌های بسیار جلب‌کننده و نامرتبط می‌توانند هنگام ارائهٔ واژه مزاحم یادگیری شوند \cite{RussoJohnson2017Touch}. به همین دلیل، نوواژه از مرحله‌های ساده، تعداد محدود گزینه و بازخورد کوتاه استفاده می‌کند.

بازی دیجیتال جای بازی واقعی را نمی‌گیرد. اثر مفید زمانی محتمل‌تر است که والد فعالیت روی صفحه را به حرکت و گفت‌وگو ادامه دهد؛ برای مثال پس از انتخاب تصویر گربه، صدای آن را تقلید کند، گربهٔ اسباب‌بازی را نشان دهد یا از کودک بخواهد آن را در کتاب پیدا کند.

\section{رسانهٔ دیجیتال و نقش والد}

رسانهٔ دیجیتال امکان پخش دقیق و تکرارشوندهٔ صوت، نمایش تصویر و ثبت پاسخ را فراهم می‌کند. محتوای آموزشی سنجیده می‌تواند فرصت یادگیری واژه ایجاد کند، اما نتیجه به کیفیت محتوا، سن کودک و زمینهٔ استفاده وابسته است \cite{Jing2023ScreenVocabulary}. کودکان خردسال همیشه آموختهٔ روی صفحه را به محیط واقعی منتقل نمی‌کنند؛ همراهی والد می‌تواند این فاصله را کاهش دهد.

برای کودکان ۱۸ تا ۲۴ ماهه، آکادمی اطفال آمریکا در صورت استفاده از رسانه، محتوای باکیفیت و استفادهٔ مشترک والد و کودک را توصیه می‌کند. برای سنین پایین‌تر از ۱۸ ماه نیز به‌جز ارتباط تصویری، استفاده از صفحه توصیه نشده است \cite{AAPMedia2016}. سازمان جهانی بهداشت نیز بر کاهش زمان نشستهٔ مبتنی بر صفحه و جایگزینی آن با حرکت، بازی و تعامل تأکید دارد \cite{WHO2019ScreenTime}.

بنابراین گروه «زیر سه سال» یک گروه یکسان نیست. برای کودک بسیار کم‌سن، تعامل مستقیم، آواز، گفت‌وگو و بازی واقعی باید فعالیت اصلی باشد. استفاده در سنین بالاتر نیز بهتر است کوتاه، همراه والد و متناسب با توانایی کودک باشد. صفحه باید آغازگر تعامل باشد، نه مقصد نهایی آن.

در طراحی نوواژه اصول زیر مورد توجه قرار گرفته‌اند:

\begin{itemize}
\item هر صفحه یک هدف اصلی و تعداد کمی گزینه دارد؛
\item ناحیه‌های لمس بزرگ و مناسب دست کودک‌اند؛
\item تصویر و صوت پیام اصلی را منتقل می‌کنند؛
\item حرکت و بازخورد به هدف آموزشی محدود می‌شوند؛
\item زمان، محتوا و خروج از حالت کودک در کنترل والد باقی می‌مانند؛
\item تبلیغ و خرید درون‌برنامه‌ای در محیط کودک وجود ندارد.
\end{itemize}

\section{کاربرد مفاهیم در نوواژه}

نوواژه هر واژه را با یک تصویر و نمونهٔ صوتی ارائه می‌کند. این ارائهٔ دیداری و شنیداری به کودک کمک می‌کند شکل شنیداری واژه را با مرجع آن مرتبط سازد. در بازی، کودک پس از شنیدن واژه تصویر مناسب را انتخاب می‌کند؛ در نتیجه فعالیت بیشتر بر شناخت دریافتی و بازیابی فعال تمرکز دارد، نه سنجش مستقیم تلفظ.

واژه‌ها در گروه‌هایی مانند حیوانات، میوه‌ها و وسایل نقلیه سازمان‌دهی شده‌اند. دسته‌بندی پیمایش محتوا را برای والد آسان می‌کند و امکان تمرین تمایز میان اعضای یک حوزهٔ معنایی را فراهم می‌آورد. در مرحله‌های آغازین می‌توان گزینه‌های متفاوت‌تری نمایش داد و با افزایش توانایی کودک، گزینه‌های نزدیک‌تر را به کار برد.

سامانه واژه‌های کمتر تمرین‌شده یا آموخته‌نشده را در اولویت قرار می‌دهد و پاسخ‌های درست را ثبت می‌کند. وضعیت «آموخته‌شده» در برنامه یک قاعدهٔ عملی برای تنظیم تمرین است و معیار بالینی محسوب نمی‌شود؛ زیرا استفادهٔ خودانگیخته، تنوع موقعیت، زمان واکنش و تعمیم به محیط واقعی را به‌طور کامل اندازه‌گیری نمی‌کند.

محتوای شخصی‌سازی‌شده یکی دیگر از قابلیت‌های برنامه است. والد می‌تواند تصویر و صدای مربوط به اعضای خانواده، اسباب‌بازی محبوب یا اشیای خانه را اضافه کند. واژه‌های مرتبط با زندگی روزمره فرصت بیشتری برای تکرار طبیعی دارند. تصویر باید واضح، صدای ضبط‌شده طبیعی و محیط ضبط کم‌صدا باشد و پس از تمرین، همان واژه در موقعیت واقعی استفاده شود.

بخش پیشرفت، تعداد تلاش‌ها، پاسخ‌های درست و وضعیت واژه‌ها را به والد نشان می‌دهد. این داده‌ها برای مشاهدهٔ روند و انتخاب تمرین بعدی مفیدند، اما آزمون استاندارد رشد زبان نیستند و نباید برای مقایسهٔ کودک با دیگران یا تشخیص اختلال استفاده شوند.

از دید فنی، برنامه‌های اندروید و آی‌اواس از منطق مشترک کاتلین چندسکویی استفاده می‌کنند و از طریق رابط \lr{API} با کارساز ارتباط دارند \cite{KotlinMultiplatformDocs2026,ComposeMultiplatformDocs2026,KtorDocs2026}. جزئیات این معماری و مراحل راه‌اندازی در فصل راهنمای فنی ارائه می‌شوند. اطلاعات حساب والد، نمایهٔ کودک، محتوای شخصی و سابقهٔ تمرین نیز باید با اصل کمینه‌سازی داده، کنترل دسترسی و ارتباط امن نگهداری شوند.

\section{جمع‌بندی}

گفتار، زبان و ارتباط مفاهیمی متمایزند و نوواژه عمدتاً بر واژگان دریافتی، پیوند صوت و تصویر، تمرین کوتاه و مشارکت والد تمرکز دارد. رشد زبان در سه سال نخست سریع اما دارای تفاوت فردی است و نگرانی‌های رشدی باید توسط متخصص بررسی شوند. گفتاردرمانی فرایندی تخصصی و فردمحور است؛ نوواژه تنها یک ابزار آموزشی مکمل برای ایجاد فرصت‌های منظم بازی و گفت‌وگو در خانواده است.

بازی می‌تواند توجه، تکرار و بازیابی فعال را تقویت کند، ولی فعالیت دیجیتال باید کوتاه، ساده و همراه والد باشد و به اشیا و تجربه‌های واقعی ادامه پیدا کند. قابلیت‌های صوت، تصویر، شخصی‌سازی و پایش پیشرفت در نوواژه بر همین اصول استوارند و نتایج آن‌ها نباید به‌عنوان ارزیابی بالینی تفسیر شوند.
